%%%%%%%%%%%%%%%%%%%%%%%%%%%%%%%%%%%%%%%%%%%%%%%%%%%%%%%%%%%%%%%%%%%%%%%%
% Plantilla TFG/TFM
% Universidad de Murcia. Facultad de Informática
% Realizado por: José Manuel Requena Plens
% Modificado: Pablo José Rocamora Zamora
% Contacto: pablojoserocamora@gmail.com
%%%%%%%%%%%%%%%%%%%%%%%%%%%%%%%%%%%%%%%%%%%%%%%%%%%%%%%%%%%%%%%%%%%%%%%%


% Lista de acrónimos (se ordenan por orden alfabético automáticamente)

% La forma de definir un acrónimo es la siguiente:
% \newacronyn{id}{siglas}{descripción}
% Donde:
% 	'id' es como vas a llamarlo desde el documento.
%	'siglas' son las siglas del acrónimo.
%	'descripción' es el texto que representan las siglas.
%
% Para usarlo en el documento tienes 4 formas:
% \gls{id} - Añade el acrónimo en su forma larga y con las siglas si es la primera vez que se utiliza, el resto de veces solo añade las siglas. (No utilices este en títulos de capítulos o secciones).
% \glsentryshort{id} - Añade solo las siglas de la id
% \glsentrylong{id} - Añade solo la descripción de la id
% \glsentryfull{id} - Añade tanto  la descripción como las siglas

\newacronym{tfg}{TFG}{Trabajo Final de Grado}
\newacronym{cryoet}{cryo-ET}{Cryo-Electron Tomography}
\newacronym{cryoem}{cryo-EM}{Cryo-Electron Microscopy}
\newacronym{tem}{TEM}{Transmission Electron Microscopy}
\newacronym{fft}{FFT}{Fast Fourier Transform}
\newacronym{gpu}{GPU}{Graphics Processing Unit}
\newacronym{cpu}{CPU}{Central Processing Unit}
\newacronym{mrc}{MRC}{Medical Research Council (formato de archivo)}
\newacronym{vtp}{VTP}{VTK PolyData}
\newacronym{sawlc}{SAWLC}{Surface-Aware Walk from Last Coordinate}
\newacronym{snr}{SNR}{Signal-to-Noise Ratio}
\newacronym{cnn}{CNN}{Convolutional Neural Network}
\newacronym{voi}{VOI}{Volume of Interest}
\newacronym{wbp}{WBP}{Weighted Back-Projection}
\newacronym{sirt}{SIRT}{Simultaneous Iterative Reconstruction Technique}
